% Macros comuns aos três ensaios de diagramação (formato-*.tex).  Cada mestre define a
% geometria e o tratamento das notas ANTES de % Macros comuns aos três ensaios de diagramação (formato-*.tex).  Cada mestre define a
% geometria e o tratamento das notas ANTES de % Macros comuns aos três ensaios de diagramação (formato-*.tex).  Cada mestre define a
% geometria e o tratamento das notas ANTES de % Macros comuns aos três ensaios de diagramação (formato-*.tex).  Cada mestre define a
% geometria e o tratamento das notas ANTES de \input{comum}; aqui ficam fonte, títulos,
% marcas de página e cabeçalhos.
\usepackage{fontspec}
\usepackage{polyglossia}
\setmainlanguage{portuguese}
\setotherlanguage{latin}
\setmainfont{EB Garamond}[Numbers={OldStyle,Proportional}]
\usepackage{microtype}
\usepackage{xcolor}
\definecolor{marca}{gray}{0.45}
\setlength{\emergencystretch}{1.5em}
\frenchspacing

% --- página do original ------------------------------------------------------
% \pg{443}: virada de página do original no corpo -> na margem externa.
% \pgn{444}: virada dentro de nota (marginpar é proibido ali) -> discreta no fluxo.
% Em modo vertical (antes de um título ou parágrafo) a marca espera o próximo
% parágrafo, senão cairia no pé da página anterior.
\newcommand{\margempg}[1]{\marginpar{\footnotesize\color{marca}[#1]}}
\gdef\pgpendente{}
\newcommand{\pg}[1]{%
  \ifvmode\xdef\pgpendente{\unexpanded\expandafter{\pgpendente}\noexpand\margempg{#1}}%
  \else\margempg{#1}\fi\ignorespaces}
\AddToHook{para/begin}{\ifinner\else
  \let\pgagora\pgpendente\gdef\pgpendente{}\pgagora\fi}
\newcommand{\pgn}[1]{{\scriptsize\color{marca}[#1]}}
% (classe book, twoside: \marginpar vai para a margem externa)
\setlength{\marginparsep}{3mm}

% --- estrutura do autor ------------------------------------------------------
% \capitulo{rótulo}{título}{summario}{pág.}
\newcommand{\capitulo}[4]{%
  \clearpage\thispagestyle{plain}\pg{#4}%
  \vspace*{2\baselineskip}%
  \begin{center}
    {\scshape\lowercase{#1}\par}\medskip
    {\large\scshape\lowercase{#2}\par}
  \end{center}
  \medskip
  {\small\noindent\textsc{summario}: --- #3\par}
  \markboth{\lowercase{#2}}{\lowercase{#2}}%
  \bigskip}
% \artigo{rótulo}{título}
\newcommand{\artigo}[2]{%
  \bigskip
  \begin{center}{\scshape\lowercase{#1}.}\\[.3ex]{\itshape #2}\end{center}
  \medskip}
% tese numerada que corre com o texto / em linha própria / subdivisão a), b)
\newcommand{\tese}[2]{\par\medskip\noindent#1.\enspace\textbf{#2}\;---\;}
\newcommand{\teseB}[2]{\par\medskip\noindent#1.\enspace\textbf{#2}\par}
\newcommand{\subp}[1]{\par\noindent\hspace*{\parindent}#1)\enspace}
% trecho omitido neste ensaio
\newcommand{\omitido}{\par{\centering[\dots]\par}}

% --- cabeçalho ---------------------------------------------------------------
\usepackage{fancyhdr}
\pagestyle{fancy}
\fancyhf{}
\fancyhead[LE]{\small\oldstylenums{\thepage}\quad\scshape theodicea}
\fancyhead[RO]{\small\scshape essencia e attributos de deus\quad\oldstylenums{\thepage}}
\renewcommand{\headrulewidth}{0pt}
\fancypagestyle{plain}{\fancyhf{}\fancyfoot[C]{\small\oldstylenums{\thepage}}}
\setcounter{page}{443}
; aqui ficam fonte, títulos,
% marcas de página e cabeçalhos.
\usepackage{fontspec}
\usepackage{polyglossia}
\setmainlanguage{portuguese}
\setotherlanguage{latin}
\setmainfont{EB Garamond}[Numbers={OldStyle,Proportional}]
\usepackage{microtype}
\usepackage{xcolor}
\definecolor{marca}{gray}{0.45}
\setlength{\emergencystretch}{1.5em}
\frenchspacing

% --- página do original ------------------------------------------------------
% \pg{443}: virada de página do original no corpo -> na margem externa.
% \pgn{444}: virada dentro de nota (marginpar é proibido ali) -> discreta no fluxo.
% Em modo vertical (antes de um título ou parágrafo) a marca espera o próximo
% parágrafo, senão cairia no pé da página anterior.
\newcommand{\margempg}[1]{\marginpar{\footnotesize\color{marca}[#1]}}
\gdef\pgpendente{}
\newcommand{\pg}[1]{%
  \ifvmode\xdef\pgpendente{\unexpanded\expandafter{\pgpendente}\noexpand\margempg{#1}}%
  \else\margempg{#1}\fi\ignorespaces}
\AddToHook{para/begin}{\ifinner\else
  \let\pgagora\pgpendente\gdef\pgpendente{}\pgagora\fi}
\newcommand{\pgn}[1]{{\scriptsize\color{marca}[#1]}}
% (classe book, twoside: \marginpar vai para a margem externa)
\setlength{\marginparsep}{3mm}

% --- estrutura do autor ------------------------------------------------------
% \capitulo{rótulo}{título}{summario}{pág.}
\newcommand{\capitulo}[4]{%
  \clearpage\thispagestyle{plain}\pg{#4}%
  \vspace*{2\baselineskip}%
  \begin{center}
    {\scshape\lowercase{#1}\par}\medskip
    {\large\scshape\lowercase{#2}\par}
  \end{center}
  \medskip
  {\small\noindent\textsc{summario}: --- #3\par}
  \markboth{\lowercase{#2}}{\lowercase{#2}}%
  \bigskip}
% \artigo{rótulo}{título}
\newcommand{\artigo}[2]{%
  \bigskip
  \begin{center}{\scshape\lowercase{#1}.}\\[.3ex]{\itshape #2}\end{center}
  \medskip}
% tese numerada que corre com o texto / em linha própria / subdivisão a), b)
\newcommand{\tese}[2]{\par\medskip\noindent#1.\enspace\textbf{#2}\;---\;}
\newcommand{\teseB}[2]{\par\medskip\noindent#1.\enspace\textbf{#2}\par}
\newcommand{\subp}[1]{\par\noindent\hspace*{\parindent}#1)\enspace}
% trecho omitido neste ensaio
\newcommand{\omitido}{\par{\centering[\dots]\par}}

% --- cabeçalho ---------------------------------------------------------------
\usepackage{fancyhdr}
\pagestyle{fancy}
\fancyhf{}
\fancyhead[LE]{\small\oldstylenums{\thepage}\quad\scshape theodicea}
\fancyhead[RO]{\small\scshape essencia e attributos de deus\quad\oldstylenums{\thepage}}
\renewcommand{\headrulewidth}{0pt}
\fancypagestyle{plain}{\fancyhf{}\fancyfoot[C]{\small\oldstylenums{\thepage}}}
\setcounter{page}{443}
; aqui ficam fonte, títulos,
% marcas de página e cabeçalhos.
\usepackage{fontspec}
\usepackage{polyglossia}
\setmainlanguage{portuguese}
\setotherlanguage{latin}
\setmainfont{EB Garamond}[Numbers={OldStyle,Proportional}]
\usepackage{microtype}
\usepackage{xcolor}
\definecolor{marca}{gray}{0.45}
\setlength{\emergencystretch}{1.5em}
\frenchspacing

% --- página do original ------------------------------------------------------
% \pg{443}: virada de página do original no corpo -> na margem externa.
% \pgn{444}: virada dentro de nota (marginpar é proibido ali) -> discreta no fluxo.
% Em modo vertical (antes de um título ou parágrafo) a marca espera o próximo
% parágrafo, senão cairia no pé da página anterior.
\newcommand{\margempg}[1]{\marginpar{\footnotesize\color{marca}[#1]}}
\gdef\pgpendente{}
\newcommand{\pg}[1]{%
  \ifvmode\xdef\pgpendente{\unexpanded\expandafter{\pgpendente}\noexpand\margempg{#1}}%
  \else\margempg{#1}\fi\ignorespaces}
\AddToHook{para/begin}{\ifinner\else
  \let\pgagora\pgpendente\gdef\pgpendente{}\pgagora\fi}
\newcommand{\pgn}[1]{{\scriptsize\color{marca}[#1]}}
% (classe book, twoside: \marginpar vai para a margem externa)
\setlength{\marginparsep}{3mm}

% --- estrutura do autor ------------------------------------------------------
% \capitulo{rótulo}{título}{summario}{pág.}
\newcommand{\capitulo}[4]{%
  \clearpage\thispagestyle{plain}\pg{#4}%
  \vspace*{2\baselineskip}%
  \begin{center}
    {\scshape\lowercase{#1}\par}\medskip
    {\large\scshape\lowercase{#2}\par}
  \end{center}
  \medskip
  {\small\noindent\textsc{summario}: --- #3\par}
  \markboth{\lowercase{#2}}{\lowercase{#2}}%
  \bigskip}
% \artigo{rótulo}{título}
\newcommand{\artigo}[2]{%
  \bigskip
  \begin{center}{\scshape\lowercase{#1}.}\\[.3ex]{\itshape #2}\end{center}
  \medskip}
% tese numerada que corre com o texto / em linha própria / subdivisão a), b)
\newcommand{\tese}[2]{\par\medskip\noindent#1.\enspace\textbf{#2}\;---\;}
\newcommand{\teseB}[2]{\par\medskip\noindent#1.\enspace\textbf{#2}\par}
\newcommand{\subp}[1]{\par\noindent\hspace*{\parindent}#1)\enspace}
% trecho omitido neste ensaio
\newcommand{\omitido}{\par{\centering[\dots]\par}}

% --- cabeçalho ---------------------------------------------------------------
\usepackage{fancyhdr}
\pagestyle{fancy}
\fancyhf{}
\fancyhead[LE]{\small\oldstylenums{\thepage}\quad\scshape theodicea}
\fancyhead[RO]{\small\scshape essencia e attributos de deus\quad\oldstylenums{\thepage}}
\renewcommand{\headrulewidth}{0pt}
\fancypagestyle{plain}{\fancyhf{}\fancyfoot[C]{\small\oldstylenums{\thepage}}}
\setcounter{page}{443}
; aqui ficam fonte, títulos,
% marcas de página e cabeçalhos.
\usepackage{fontspec}
\usepackage{polyglossia}
\setmainlanguage{portuguese}
\setotherlanguage{latin}
\setmainfont{EB Garamond}[Numbers={OldStyle,Proportional}]
\usepackage{microtype}
\usepackage{xcolor}
\definecolor{marca}{gray}{0.45}
\setlength{\emergencystretch}{1.5em}
\frenchspacing

% --- página do original ------------------------------------------------------
% \pg{443}: virada de página do original no corpo -> na margem externa.
% \pgn{444}: virada dentro de nota (marginpar é proibido ali) -> discreta no fluxo.
% Em modo vertical (antes de um título ou parágrafo) a marca espera o próximo
% parágrafo, senão cairia no pé da página anterior.
\newcommand{\margempg}[1]{\marginpar{\footnotesize\color{marca}[#1]}}
\gdef\pgpendente{}
\newcommand{\pg}[1]{%
  \ifvmode\xdef\pgpendente{\unexpanded\expandafter{\pgpendente}\noexpand\margempg{#1}}%
  \else\margempg{#1}\fi\ignorespaces}
\AddToHook{para/begin}{\ifinner\else
  \let\pgagora\pgpendente\gdef\pgpendente{}\pgagora\fi}
\newcommand{\pgn}[1]{{\scriptsize\color{marca}[#1]}}
% (classe book, twoside: \marginpar vai para a margem externa)
\setlength{\marginparsep}{3mm}

% --- estrutura do autor ------------------------------------------------------
% \capitulo{rótulo}{título}{summario}{pág.}
\newcommand{\capitulo}[4]{%
  \clearpage\thispagestyle{plain}\pg{#4}%
  \vspace*{2\baselineskip}%
  \begin{center}
    {\scshape\lowercase{#1}\par}\medskip
    {\large\scshape\lowercase{#2}\par}
  \end{center}
  \medskip
  {\small\noindent\textsc{summario}: --- #3\par}
  \markboth{\lowercase{#2}}{\lowercase{#2}}%
  \bigskip}
% \artigo{rótulo}{título}
\newcommand{\artigo}[2]{%
  \bigskip
  \begin{center}{\scshape\lowercase{#1}.}\\[.3ex]{\itshape #2}\end{center}
  \medskip}
% tese numerada que corre com o texto / em linha própria / subdivisão a), b)
\newcommand{\tese}[2]{\par\medskip\noindent#1.\enspace\textbf{#2}\;---\;}
\newcommand{\teseB}[2]{\par\medskip\noindent#1.\enspace\textbf{#2}\par}
\newcommand{\subp}[1]{\par\noindent\hspace*{\parindent}#1)\enspace}
% trecho omitido neste ensaio
\newcommand{\omitido}{\par{\centering[\dots]\par}}

% --- cabeçalho ---------------------------------------------------------------
\usepackage{fancyhdr}
\pagestyle{fancy}
\fancyhf{}
\fancyhead[LE]{\small\oldstylenums{\thepage}\quad\scshape theodicea}
\fancyhead[RO]{\small\scshape essencia e attributos de deus\quad\oldstylenums{\thepage}}
\renewcommand{\headrulewidth}{0pt}
\fancypagestyle{plain}{\fancyhf{}\fancyfoot[C]{\small\oldstylenums{\thepage}}}
\setcounter{page}{443}
